% Generated draft: requires mathematical and source review.
\documentclass[11pt]{article}
\usepackage[margin=25mm]{geometry}
\usepackage{fontspec}
\setmainfont{Latin Modern Roman}
\usepackage{amsmath,amssymb,mathtools,unicode-math}
\setmathfont{Latin Modern Math}
\usepackage{enumitem,needspace,xurl,hyperref}
\setlength{\parindent}{0pt}
\setlength{\parskip}{5pt}
\setlength{\emergencystretch}{3em}
\newcommand{\N}{\mathbb{N}}
\newcommand{\Z}{\mathbb{Z}}
\newcommand{\Q}{\mathbb{Q}}
\newcommand{\R}{\mathbb{R}}
\newcommand{\C}{\mathbb{C}}
\newcommand{\F}{\mathbb{F}}
\newcommand{\A}{\mathbb{A}}
\newcommand{\PP}{\mathbb{P}}
\newcommand{\E}{\mathbb{E}}
\newcommand{\eps}{\varepsilon}
\newcommand{\dd}{\,\mathrm{d}}
\begin{document}
\section*{UnsolvedMath ID: 30005422}
\textbf{Source problem code:} catalogue entry
\subsection{Definitions and mathematical statement}
The supplied record refers to a mutation theorem for \(0\)-Auslander extriangulated categories, identified in the original excerpt only as ``Theorem 2.'' Neither that theorem nor definitions of its hypotheses, the objects undergoing mutation, or the mutation operation are supplied. Consequently, a complete mathematical formulation with explicit domains and quantifiers cannot be recovered from this record alone.

\textbf{Recorded mathematical target.} Determine whether the referenced Theorem 2 can be applied to discover new mutations. The supplied cleaned statement additionally asks which established mutation theories can be recovered from the mutation theorem for \(0\)-Auslander extriangulated categories, and which new mutation operations arise from it. This is a research question whose precise formulation requires the missing theorem and mutation data; it does not assert that any new mutation exists. The meaning of recovery, the comparison class of established theories, and the criterion for novelty are unspecified.

\textbf{Qualifications from the original excerpt.} The source emphasizes that its examples have different categorical structures: \((\mathcal C,\Delta_T)\) is a relative structure of a triangulated category; \(K^{[-1,0]}(\operatorname{proj}\Lambda)\) is extension-closed in a triangulated category; and \(\mathcal W\) is an exact category whose quotient by the ideal of morphisms factoring through some projective-injective object is, in the example described, neither exact nor triangulated. That quotient is described as relevant from the combinatorial perspective. The excerpt presents extriangulated structures as a framework accommodating these different situations. These are qualifications concerning the referenced examples, not universal assertions about such categories or quotients. The record does not specify \(\mathcal C\), \(T\), \(\Delta_T\), \(\Lambda\), \(\mathcal W\), the relevant projective-injective objects, or the hypotheses connecting these examples to Theorem 2.
\subsection{Short English statement}
The question asks whether a referenced mutation theorem produces new mutations; the cleaned record also asks which existing mutation theories it recovers. The theorem, its hypotheses, and the mutation data are missing, so the supplied material supports only this qualified research target. The different categorical settings in the excerpt must be retained when a precise statement is reconstructed.
\subsection{Sources}
\begin{itemize}
\item Supplied local UnsolvedMath registry record, permanent ID \(30005422\), exact source problem code \texttt{catalogue entry}, title ``Mutation Theories from Extriangulated Categories,'' in the Oberwolfach Reports Open Problems collection. No individual proposer is supplied; the recorded proposed year is 2023.
\item \emph{Representation Theory of Quivers and Finite-Dimensional Algebras} (2023), Oberwolfach Reports, pp.~397--486, DOI \texttt{10.4171/owr/2023/7}. Supplied source and literature link: \url{https://doi.org/10.4171/owr/2023/7}. The extracted text includes a page marker for p.~415. The report is supplied as the authoritative source, but the referenced Theorem 2 is absent from the record.
\item Supplied literature assessment dated \texttt{2026-08-22}: the record labels the problem open and states that the extracted question lacks theorem and mutation data, and that its curation located no authoritative source resolving this under-specified problem. The listed source checked is the ``original OWR report,'' \url{https://doi.org/10.4171/owr/2023/7}. This is supplied dated curation, not an independently verified literature assessment.
\item The registry assigns the statement status \texttt{corrected\_verified} and reports that the cleaned statement was checked against the official report and preserves its hypotheses, notation, and target. This is a supplied verification claim, not verification performed for this draft; the missing theorem prevents checking that claim from the supplied excerpt.
\end{itemize}
\end{document}
